\section{Limitations and future work}\label{sec:conclusion}

\paragraph{The relational analysis has not met a real secret} None of the 30 real
workflows declares or implies a secret in its source's threat model, so the analysis
that motivates this paper verified nothing on them. The pattern it targets, a branch on
private data that is not itself sent to a model (a customer's tier, a risk flag, the
verdict of a local detector of personal data), is plausible and easy to write, but we
have not shown that deployed workflows contain it. Finding such workflows is the most
important next step.

\paragraph{Declassification is too coarse} Real workflows send private data to a
provider, which in OrchLang requires declassifying it, which makes it public to every
observer. A policy relative to observers, under which the provider may read content
that the network and the bill may not reveal, would let the analysis reason about those
workflows. \Cref{thm:ni,thm:leakage} would then need an equivalence of histories indexed
by observer.

\paragraph{Integrity over-reports} A value computed in an arm guarded by untrusted data
is labelled untrusted even when it is computed only from trusted data. Separating data
and program-counter labels for integrity, as \cref{sec:labels} does for confidentiality,
would remove the two false positives of the held-out split.

\paragraph{Retry cannot carry feedback} Self-correction loops that show the model its
earlier attempts had to be unrolled in five of the thirty ports. A retry whose body can
read an accumulator of declared byte growth would keep the bound finite and express
these loops directly.

\paragraph{Agents are out of scope by design} Thirteen of the 22 exclusions are workflows
whose shape is decided at run time, nine of them because the model chooses what runs
next. OrchLang's fixed shape is what makes its analyses decidable; a type system for
agent loops would need bounds on the model's choices.

\paragraph{The guaranteed bound is loose} On the real workflows it is four times the
estimate where requests carry only inputs and up to 128 times where they carry earlier
responses, unless models declare client-side byte caps.

\paragraph{The certificate is not independently checked, and the implementation is not
verified} A small checker that re-derives the bounds from the certificate and the source,
tested against tampered certificates, would turn the certificate into evidence. The C++
implementation is tested against the Coq model through the runtime, not verified against
it.

\section{Conclusion}\label{sec:conclusion-final}

An analysis that compares LLM requests by anything coarser than their content cannot be
sound against providers that answer content, unless it gives up on every secret branch.
Comparing content gives noninterference that needs no assumption about providers,
tokenizers or observers, and a leakage bound that counts request classes rather than
noisy observations. Token bounds have to be built from the one measure that adds up,
bytes, and converted through tokenizer contracts that measurements can refute. Both
results are mechanised for a core calculus, implemented in a compiler, and tested by a
harness that can fail. On 30 real workflows the bounds held and the flow checks found
the places where untrusted content decides an effect, but the side channel we set out
to close did not appear. Whether it matters in practice depends on whether deployed
workflows branch on private data that they do not send to a model, and that question
remains open.
